problem:  
**Problem (Boundary \(K\)-Invariant for Non‑Globally Hypoelliptic Hörmander Operators).**  
Let \( (M,\mathcal H) \) be a compact manifold with a bracket‑generating distribution \(\mathcal H\) satisfying Hörmander’s condition, and let \(P\) be a Rockland differential operator on \((M,\mathcal H)\) with principal symbol \(\sigma_H(P)\in C^*(T_{\mathcal H}M)\).  
Consider the **filtered tangent groupoid** short exact sequence of \(C^*\)-algebras  
\[
0 \;\longrightarrow\; J \;\longrightarrow\; C^*(\mathbb T_{\mathcal H}M) \;\xrightarrow{\pi_0}\; C^*(T_{\mathcal H}M) \;\longrightarrow\; 0,
\]
where \(J\) is generated by the fibers for \(t>0\) (the analysis at finite scale).

This induces the boundary map in \(K\)-theory,
\[
\partial: K_0(C^*(T_{\mathcal H}M))\longrightarrow K_1(J),
\]
which measures the failure of a symbol class to extend to an invertible analytic family through the deformation.

---

**Conjectural problem:**  
1. (*Existence*) For each Rockland operator \(P\), define  
\[
\Xi(P) := \partial([\sigma_H(P)]) \in K_1(J).
\]  
Does \(\Xi(P)\) detect **global hypoellipticity** in the sense that  
\[
\Xi(P)=0 \quad\Longleftrightarrow\quad P\text{ is globally hypoelliptic on }M?
\]

2. (*Analytic characterization*) Can \(\Xi(P)\) be expressed as the spectral‐flow class of the family of resolvents  
\((P_t+i)^{-1}\), \(t>0\),  
in the continuous field of operator algebras \(C^*(\mathbb T_{\mathcal H}M|_{t>0})\)?  

3. (*Geometric dependence*) Under smooth deformations of the Hörmander structure \(\mathcal H_t\) preserving the local nilpotent type, does \(\Xi(P_t)\) remain constant in \(K_1(J)\)? If not, can its variation be related to the formation or loss of global hypoellipticity within the family?

---

**Why it is a "good" problem:**  
This modified problem removes the vagueness of the earlier proposals by grounding the *secondary invariant* in an exact and canonical analytic construction: the **boundary map in \(K\)-theory** for the filtered tangent groupoid. It transforms the elusive property of *global* hypoellipticity into a precise question about the vanishing of a boundary class \(\Xi(P)\).  

If \(\Xi(P)\) indeed vanishes exactly for globally hypoelliptic operators, it would unveil a new quantitative bridge between topological \(K\)-theory and global analytic regularity.  Conversely, non‑vanishing of \(\Xi(P)\) would represent a measurable obstruction—a secondary invariant complementing the Rockland symbol condition.  

The problem is simple to state yet conceptually rich, demanding the fusion of microlocal analysis, operator algebra exact sequences, and spectral‑flow techniques.  It thus fits perfectly the profile of a **“good” research problem**: precise, foundational, and with the potential to reveal an entirely new layer of structure in the interplay between Hörmander hypoellipticity and noncommutative geometry.